% Master's thesis (work-in-progress draft). Compile with pdfLaTeX + BibTeX, e.g. on Overleaf.
% One-sided layout for the digital draft: no blank pages before chapters and the same margins on
% every page. For a printed copy, use twoside,binding=0.6cm.
\documentclass[a4paper,oneside,english]{sapthesis}

\usepackage[english]{babel}
\usepackage[utf8]{inputenc}
\usepackage{microtype}
\usepackage{amsmath,amssymb}
\usepackage{booktabs}
\usepackage{tabularx}
\usepackage{graphicx}
\usepackage{siunitx}
\usepackage[numbers,sort&compress]{natbib}
\usepackage{hyperref}
\hypersetup{pdftitle={Conformal Triage for Thyroid Nodules with a Radiomics-Prompted Language Model},
            pdfauthor={Sezer Mezgil}}

% TODO: working title, to be fixed with the advisor.
\title{Conformal Triage for Thyroid Nodules\\with a Radiomics-Prompted Language Model}
\author{Sezer Mezgil}
\IDnumber{2180600}
\course{Laurea Magistrale in Data Science}
\courseorganizer{Facoltà di Ingegneria dell'Informazione, Informatica e Statistica}
\AcademicYear{2025/2026} % TODO: confirm the academic year of graduation
\advisor{Prof. Pierpaolo Brutti}
\authoremail{mezgil.2180600@studenti.uniroma1.it}
\copyyear{2026}
\thesistype{Master's thesis}
\versiondate{Draft, \today}

\begin{document}

\frontmatter
\maketitle

\begin{abstract}
Machine-learning models for thyroid ultrasound are rarely accompanied by a guarantee on their
errors. This thesis studies a triage layer that decides automatically only where the error rate can
be certified: at most 5\% cancers among nodules called benign and at most 20\% benign nodules among
those called malignant, with all other nodules referred to a clinician. The thresholds are set with
Learn then Test, a conformal risk-control method. The underlying classifier follows Ra et al.\
(2025): radiomics features of the nodule are written as text and read by a language model
fine-tuned with QLoRA.

On the TN3K dataset (3,493 images with biopsy-based labels), an audit showed that the official
calibration and test sources are not exchangeable and that the original features had defects; a
pooled split and a revised extraction address both. Six features computed from TI-RADS signs,
such as the taller-than-wide ratio and the solidity of the outline, give the largest gain: over
three matched training seeds they raise the test AUC of Qwen2.5-1.5B from 0.726 to 0.785
($\pm 0.018$). XGBoost on the same features reaches 0.835, and combining the two models adds
nothing to XGBoost. Larger backbones (Qwen2.5-7B, Llama-2-7B with the settings of Ra et al.) and
alternative prompt formats did not help in the settings tried. The planned Learn-then-Test procedure certifies no
automatic decision, for the language model or for XGBoost: the cleanest low-risk groups are too
small or not clean enough for the evidence required, and the planned order of tests stops at
thresholds that hold too few nodules. A design fixed in advance, with thresholds at quantiles of
the validation scores, certifies the automatic malignant calls in 99\% of 2,000 re-splits of the
data; the automatic benign calls remain out of reach. Guarantees on the share of missed cancers, from split
conformal prediction and a variant of Learn then Test, are attainable and hold over 4,000 re-splits
of the data, but about 10\% of the automatic benign calls are still cancers (6\% with XGBoost).
Which guarantee is attainable, and what it is worth clinically, are the main results.
\end{abstract}

\tableofcontents

\mainmatter
\chapter{Introduction}
\label{ch:introduction}

\section{Motivation}

Most thyroid nodules are benign, but deciding which ones need a biopsy is a daily clinical task.
Ultrasound risk systems guide this decision~\citep{tessler2017}, and fine-needle aspiration
settles it. Machine-learning models can estimate the risk of malignancy from an ultrasound image,
but a model that is right most of the time is of limited use if nobody knows which of its
decisions can be trusted.

This thesis studies a model that decides only where it can guarantee its error rate. Nodules it
calls benign must contain at most 5\% cancers, nodules it calls malignant at most 20\% benign ones,
and every other nodule is referred to a clinician. The guarantee is obtained with Learn then
Test~\citep{angelopoulos2021ltt}, a conformal risk-control method that needs no assumption on the
model, only calibration data from the same distribution as future cases.

The underlying classifier follows Ra et al.~\citep{ra2025}: radiomics features are extracted from
the nodule, written as text, and read by a fine-tuned language model. The approach is attractive
because a language model could use what feature names mean, and because its answers can be
queried repeatedly to measure how stable they are.

\section{Research questions}

\begin{enumerate}
  \item[RQ1] How well does a language model classify thyroid nodules from radiomics prompts, and
             how does it compare with a classical model on the same features?
  \item[RQ2] Which choices of the pipeline matter: the feature extraction, the way values are
             written, and the size of the model?
  \item[RQ3] Can Learn then Test certify automatic decisions at clinically meaningful error rates,
             and does agreement across perturbed queries help?
\end{enumerate}

\section{Contributions}

\begin{itemize}
  \item An adaptation of the radiomics-to-language-model pipeline of Ra et al.\ to thyroid
        ultrasound (TN3K), with a three-way triage layer calibrated by Learn then Test.
  \item A data audit of TN3K: the official calibration and test sources are not exchangeable, and
        the original feature extraction had three defects (few grey levels, size in pixels,
        multi-component masks). A pooled split and a revised extraction address them.
  \item Six features computed from TI-RADS signs that are missing from standard radiomics. They
        raise the language model's test AUC by about 0.06 over three matched seeds, more than any
        other change, and XGBoost's by 0.026.
  \item A controlled comparison with a tabular model: with the same features the language model
        stays about 0.05 AUC below XGBoost, a fusion of the two shows it adds no information, and
        neither larger backbones, the settings of Ra et al.\ nor plain-language prompts change
        this.
  \item A negative result with its explanation: no automatic decision could be certified at the
        planned error rates, for the language model or for XGBoost, because the cleanest low-risk
        groups are too small or not clean enough, and because the planned order of tests stops at
        near-empty thresholds. Perturbation votes do not beat the softmax probability.
  \item A map of what is attainable instead: guarantees on missed cancers hold, verified over
        4,000 re-splits, but leave about 10\% cancers among the automatic benign calls, a risk at
        which biopsy is still recommended.
\end{itemize}

\section{Structure of the thesis}

Chapter~\ref{ch:background} introduces the clinical and methodological background.
Chapter~\ref{ch:data} describes the data, its quality issues and the splits.
Chapter~\ref{ch:methods} presents the pipeline and the triage layer.
Chapter~\ref{ch:results} reports the experiments, organised by question, and
Chapter~\ref{ch:discussion} discusses them. Chapter~\ref{ch:future} lists the ongoing and planned
work.

\chapter{Background and related work}
\label{ch:background}

\section{Thyroid nodules, TI-RADS and the reference standard}
\label{sec:bg-thyroid}

Thyroid nodules are common and most of them are benign. Ultrasound is the first examination, and
risk systems such as ACR TI-RADS~\citep{tessler2017} score a nodule from its appearance:
composition, echogenicity, shape (in particular a taller-than-wide shape), margin and echogenic
foci. The score and the size of the nodule decide whether a fine-needle aspiration (FNA) is
recommended. The cytology of the aspirate, and surgery where needed, establish the diagnosis.

FNA is the reference standard but it is invasive and uses clinical resources, and many aspirated
nodules turn out benign. A system that could clear clearly benign nodules with a known error
rate, and pass all doubtful ones to a clinician, would reduce this load without hiding its
mistakes.

\section{Radiomics}
\label{sec:bg-radiomics}

Radiomics describes a region of an image with many hand-crafted, quantitative features: shape
(size, elongation, roundness, outline), first-order statistics of the intensities inside the
region, and texture features computed from matrices that count how grey levels are arranged
(co-occurrence, run lengths, size zones, dependences, neighbourhood differences). PyRadiomics is a
widely used open implementation~\citep{vangriethuysen2017}. Texture features depend on how the
intensities are discretised into grey levels, and shape features depend on the physical scale of
the image; both matter for TN3K (Section~\ref{sec:methods-extraction}).

Radiomics yields many correlated features for a modest number of cases, so a feature selection
step is usual. Minimum-redundancy maximum-relevance (mRMR)~\citep{peng2005} adds features one by one,
preferring those related to the label and unrelated to the features already chosen.

\section{Language models for radiomics and tabular data}
\label{sec:bg-llm}

A row of tabular data can be written as text, ``feature name: value'', and given to a pre-trained
language model. The idea is that the model may link feature names to knowledge acquired in
pre-training, which a classical model that sees only a column of numbers cannot do. Fine-tuning
all weights of a large model is expensive; low-rank adaptation (LoRA)~\citep{hu2022lora} trains
small added matrices instead, and QLoRA~\citep{dettmers2023qlora} does so on a 4-bit copy of the
frozen model, which makes 7B-parameter models trainable on a single mid-range GPU.

\subsection{Ra et al.\ (2025)}

Ra et al.~\citep{ra2025} apply this idea to mammography. They extract 67 PyRadiomics features
(shape, first-order, GLCM, GLSZM) from segmented masses, select 8 with mRMR, write them as
sentences, and fine-tune Llama-2-7B with LoRA (rank 16) and a classification head. On VinDr-Mammo
(235 test images) their model reaches AUC 0.700 against 0.605--0.727 for classical and tabular
baselines; on INbreast, after retraining on half of the data, 0.862. They also report that
replacing the feature names with generic labels lowers performance, and that their model can be
extended to a new feature set without retraining from scratch. This thesis adapts their pipeline
to thyroid ultrasound; a critical reading of their evaluation is given in
Section~\ref{sec:disc-llm}.

\section{Conformal prediction and risk control}
\label{sec:bg-conformal}

Conformal prediction~\citep{vovk2005} turns any model score into prediction sets with a
finite-sample guarantee: the true label lies in the set with probability at least $1 - \alpha$,
provided the calibration data and the new data are exchangeable. It makes no assumption on the
model or the data distribution. With two classes, a set is either one label (a decision), both
labels (``unsure'') or empty, and its guarantee concerns coverage, not the error rate among the
decisions. Class-conditional (Mondrian) versions give the guarantee separately for each class.

Selective classification lets a model abstain and bounds the error rate among the cases it does
decide~\citep{geifman2017}. Learn then Test (LTT)~\citep{angelopoulos2021ltt} generalises this: any
set of risks, written as expected losses, can be controlled by testing candidate parameters as
hypotheses on calibration data, with a guarantee that holds with probability $1 - \delta$. Conformal
risk control~\citep{angelopoulos2024crc} bounds the expected value of a monotone loss, such as the
share of positives that are missed. All of these assume exchangeability. When the test data are
shifted, weighted versions can restore the guarantee under covariate shift~\citep{tibshirani2019}.

\section{Uncertainty from repeated queries}
\label{sec:bg-queries}

For language models, a common uncertainty signal is agreement across repeated or perturbed
queries: self-consistency over sampled answers~\citep{wang2023selfconsistency}, semantic entropy
over answers with the same meaning~\citep{kuhn2023semantic}, and conformal prediction built on
sampling frequencies when the model's probabilities are not accessible~\citep{su2024api}. A recent
scoping review~\citep{ashby2026} finds that such sampling-based signals often outperform those
based on the model's probabilities, mostly on general language benchmarks. Noorani et
al.~\citep{noorani2025} study how many queries are needed for a conformal guarantee. Whether these
findings carry over to a fine-tuned classifier reading radiomics prompts is one of the questions of
this thesis.

\chapter{Data}
\label{ch:data}

\section{The TN3K dataset}
% 3,493 ultrasound images from 2,421 patients, one nodule mask per image, benign/malignant label
% from cytological biopsy \citep{gong2021tn3k,gong2022acl}. Official patient-disjoint split:
% trainval 2,879, test 614.

\section{Image and mask quality control}
% Report §11.4: image sizes differ between the official sources; multi-component and
% border-touching masks; about 9 grey levels inside a nodule with binWidth 25.

\section{The official split is not exchangeable}
% Report §5: two-sample AUC 0.63 between calibration and official test; image-size artefact;
% sources differ beyond pixel scale (A5 check in runs/2026-10-08_v2_pipeline).

\section{The pooled split}
% splits/pooled_seed42.csv: stratified by label and source; train 2,095, validation 436,
% calibration 524, test 438. Exchangeability check after re-splitting.

\subsection{Patient-level independence}
% Image-level split without patient IDs: limitation text agreed on 2026-10-09.

\chapter{Methods}
\label{ch:methods}

\section{Overview of the pipeline}
\label{sec:methods-overview}

The pipeline follows Ra et al.~\citep{ra2025}, adapted from mammography to thyroid ultrasound,
and adds a triage layer on top. For each image:

\begin{enumerate}
  \item radiomics features are extracted from the image inside the nodule mask
        (Section~\ref{sec:methods-extraction});
  \item the features are standardised and a small subset is selected, using the training split
        only (Section~\ref{sec:methods-selection});
  \item the selected feature names and values are written as a short English text
        (Section~\ref{sec:methods-prompts});
  \item a pre-trained language model, fine-tuned for this task, reads the text and outputs a
        probability of malignancy (Section~\ref{sec:methods-llm});
  \item a three-way rule turns this score into \emph{auto-benign}, \emph{auto-malignant} or
        \emph{refer to a clinician}, with thresholds calibrated so that the error rates of the
        automatic decisions are statistically guaranteed (Section~\ref{sec:methods-triage}).
\end{enumerate}

An XGBoost classifier trained on the same features serves as the benchmark
(Section~\ref{sec:methods-xgb}). All code is in the thesis repository: the triage layer and
its unit tests in \texttt{conformal\_triage/} and \texttt{tests/}, the experiments in
\texttt{experiments/}, and one record per experiment in \texttt{runs/}.

\section{Radiomics feature extraction}
\label{sec:methods-extraction}

Features are extracted with PyRadiomics~\citep{vangriethuysen2017} in 2D mode, from the original
image only (no wavelet or Laplacian-of-Gaussian filtered images, as in~\citep{ra2025}). Each image
is converted to greyscale and rescaled to the range 0--255 by its own minimum and maximum, and
each mask is binarised at 127 (Section~\ref{sec:data-tn3k}). Because no pixel spacing is
available, the spacing is set to 1, so all sizes are in pixels.

\subsection{First extraction (v1)}

The first extraction used four feature families: shape (shape2D), first-order intensity
statistics, the grey-level co-occurrence matrix (GLCM) and the grey-level size-zone matrix
(GLSZM), giving 67 usable features. This is the set used by Ra et al. Texture features were
computed with a fixed bin width of 25 grey values.

The quality checks of Section~\ref{sec:data-qc} revealed three problems with this extraction:

\begin{itemize}
  \item with a bin width of 25, a nodule contains a median of about 9 grey levels (interquartile
        range 7--10), which leaves little room for texture;
  \item shape features in pixels partly measure the size of the image;
  \item multi-component masks were used as they are, so a single shape feature could describe
        several separate regions.
\end{itemize}

\subsection{Revised extraction (v2)}

The revised extraction changes four settings and keeps everything else:

\begin{itemize}
  \item \textbf{Fixed bin count.} Texture features use 64 grey levels per image instead of a fixed
        bin width.
  \item \textbf{Size relative to the image.} With image width $W$ and height $H$, lengths
        (major and minor axis, maximum diameter, perimeter) are divided by $\sqrt{WH}$, areas by
        $WH$, and the perimeter-to-surface ratio is multiplied by $\sqrt{WH}$. Shape features
        without units (sphericity, elongation) are unchanged.
  \item \textbf{Three more texture families:} the grey-level run-length matrix (GLRLM), the
        grey-level dependence matrix (GLDM) and the neighbouring grey-tone difference matrix
        (NGTDM).
  \item \textbf{Largest component.} For multi-component masks only the largest region is kept,
        a rule fixed before extraction so that each image describes one region.
\end{itemize}

This gives 102 features for all 3,493 images, with no extraction errors. As a consistency check,
on single-component masks the features shared with v1 reproduce the v1 values exactly.

\section{TI-RADS-motivated features}
\label{sec:methods-clinical}

PyRadiomics describes the nodule region on its own: its shape features do not depend on
orientation, and every feature is computed inside the mask. Several TI-RADS signs depend on
exactly what this leaves out, the orientation of the nodule and its contrast with the surrounding
tissue~\citep{tessler2017}. Following the TI-RADS categories, six features were computed from each
image as stored (without per-image rescaling) and its mask, with settings fixed before any label
was looked at (\texttt{conformal\_triage/clinical.py}):

\begin{itemize}
  \item \textbf{Shape:} the taller-than-wide ratio, the vertical extent of the mask divided by its
        horizontal extent. The image is taken as acquired, with depth downwards; whether a view is
        transverse or longitudinal is not recorded in TN3K.
  \item \textbf{Echogenicity:} the mean brightness inside the nodule divided by the mean of a
        ring of surrounding tissue (width a quarter of the nodule's equivalent radius, at least 4
        pixels, excluding the black area outside the ultrasound sector).
  \item \textbf{Echogenic foci:} small bright spots per 1,000 nodule pixels. A spot is a
        connected region of 2--40 pixels inside the slightly eroded nodule whose white top-hat
        response (radius 3) exceeds the median by six robust standard deviations and which is
        brighter than the surrounding ring.
  \item \textbf{Margin:} solidity, the nodule area divided by the area of its convex hull
        (below 1 for lobulated or irregular outlines), and margin sharpness, the mean gradient
        magnitude on the boundary divided by the brightness spread of the nodule and the ring.
        ACR TI-RADS scores an ill-defined margin with 0 points; its margin points come from a
        lobulated or irregular outline (2) or extrathyroidal extension (3). Solidity targets the
        scored sign; margin sharpness is an additional measurement inspired by the category, not
        a counterpart of a scored sign.
  \item \textbf{Composition:} the share of nodule pixels darker than a quarter of the ring's
        mean brightness, as fluid would be.
\end{itemize}

All six are ratios or densities, so they do not depend on image size or on the overall gain of
the scanner. They are heuristics: each was checked on synthetic nodules (unit tests), not against
expert TI-RADS ratings, which TN3K does not provide. When used, they are always placed first in
the prompt; the radiomics features are still filtered and ranked by mRMR as described below.

\section{Preprocessing and feature selection}
\label{sec:methods-selection}

All steps below are fitted on the training split and then applied unchanged to the other splits:

\begin{enumerate}
  \item infinite values are set to missing and constant features are removed;
  \item missing values are replaced by the training median, and every feature is standardised to
        zero mean and unit variance (z-score);
  \item \textbf{near-duplicate filter} (later experiments): of two features with absolute
        Pearson correlation above 0.95, only the one more correlated with the label is kept;
  \item \textbf{mRMR}~\citep{peng2005} (minimum redundancy, maximum relevance) ranks the remaining
        features, and the first $k$ are used: $k = 8$ in the first experiments, as in Ra et al.,
        and $k = 16$ later.
\end{enumerate}

The near-duplicate filter was added after mRMR selected two features that are the same number in
2D (pixel surface and mesh surface, both relative to the image). The mRMR implementation penalises
a candidate by its \emph{average} correlation with the features already chosen, so a copy of one
chosen feature can still be picked when it is uncorrelated with the others. The selected features
keep their mRMR order in the prompt.

\section{From features to prompts}
\label{sec:methods-prompts}

Each nodule becomes one text: a fixed instruction followed by one sentence per selected feature.
Feature names are made readable (\texttt{original\_shape2D\_Elongation} becomes
\emph{Elongation}). With the format used in all main experiments (\emph{decimal}), a prompt begins:

\begin{quote}\small
Classify this thyroid nodule as benign or malignant based on the following standardized radiomic
features. The 'Major Axis Length Relative' feature is measured at 0.335. The 'Elongation' feature
is measured at $-$1.140. \dots
\end{quote}

The value is the z-score from Section~\ref{sec:methods-selection}, written with three decimals.
Four formats were compared (Table~\ref{tab:methods-formats}); a fifth is implemented but has not
been run yet (Section~\ref{sec:future-prompts}).

\begin{table}[ht]
  \centering
  \caption{Value formats in the prompt.}
  \label{tab:methods-formats}
  \begin{tabularx}{\linewidth}{l>{\raggedright\arraybackslash}X}
    \toprule
    format & sentence for one feature \\
    \midrule
    decimal    & The '\dots' feature is measured at $-$1.140. \\
    z-score    & The '\dots' feature has a z-score of $-$1.14 relative to the training nodules. \\
    percentile & The '\dots' feature is at percentile 12 of the training nodules. \\
    \midrule
    rich & z-score, percentile and a plain-language band, e.g.\ ``lower than
                         typical, meaning a less round, more irregular outline'' \\
    semantic only (not run yet) & the band and phrase without numbers \\
    \bottomrule
  \end{tabularx}
\end{table}

\emph{Decimal} and \emph{z-score} carry the same number; they differ only in saying what the
number is. \emph{Percentile} replaces the z-score by its rank among the training nodules, which
removes the information on how extreme a value is. The phrases of the two planned formats
describe what a feature measures, never which direction is malignant.

\section{Language-model classifier}
\label{sec:methods-llm}

\paragraph{Architecture.} The language model is used as a classifier, not as a text generator. A
linear classification head with two outputs is placed on the hidden state of the last token and
initialised at random; the probability of malignancy is the softmax of its two outputs. Ra et al.\
also train a classification head, with a binary cross-entropy loss, and compare it with training
the head alone.

\paragraph{Fine-tuning.} The pre-trained weights are loaded in 4-bit precision (NF4 with double
quantisation) and kept frozen; only low-rank adapters (LoRA) and the classification head are
trained (QLoRA~\citep{hu2022lora,dettmers2023qlora}). Table~\ref{tab:methods-llm} lists the
settings.

\begin{table}[ht]
  \centering
  \small
  \caption{Fine-tuning settings, with those reported by Ra et al.~\citep{ra2025}.}
  \label{tab:methods-llm}
  \begin{tabular}{lp{6.2cm}p{3.6cm}}
    \toprule
    setting & this work & Ra et al. \\
    \midrule
    backbone        & Qwen2.5-1.5B~\citep{qwen2024}; also Qwen2.5-7B and Llama-2-7B & Llama-2-7B~\citep{touvron2023llama2} \\
    weights         & 4-bit, frozen & not stated \\
    LoRA            & rank 16, $\alpha = 32$, dropout 0.05 & rank 16 \\
    LoRA placement  & attention Q and V (all linear layers in one run) & all layers \\
    loss            & class-weighted cross-entropy & binary cross-entropy \\
    optimiser       & AdamW, learning rate $2 \times 10^{-4}$ (Llama-2-7B: $10^{-4}$; one Qwen2.5-7B run: $5 \times 10^{-5}$), cosine schedule, 5\% warm-up & AdamW, $10^{-4}$, cosine schedule \\
    batch           & 8 $\times$ 2 accumulation steps (7B: 4 $\times$ 4) & not stated \\
    epochs          & 5 & over 500 \\
    model selection & epoch with the best validation AUC & best on validation \\
    \bottomrule
  \end{tabular}
\end{table}

Class weights are $n / (2 n_c)$ for a class with $n_c$ of the $n$ training nodules, so that
errors on the less frequent malignant class cost more. The maximum prompt length is checked
before training so that no prompt is truncated.

\paragraph{Hardware and precision.} The first models were trained on a free Colab T4 GPU in
16-bit floating point (fp16), the later ones on an L4 GPU in bfloat16. Probabilities are computed
in full precision after training. Earlier runs predicted in bfloat16, which rounded the scores so
that many nodules shared the same value (195 distinct values among 524 calibration nodules); this
barely changes the AUC but coarsens the thresholds of the triage layer.

\section{Classical benchmark}
\label{sec:methods-xgb}

XGBoost~\citep{chen2016xgboost} is trained on the same split, with the same preprocessing and the
same mRMR order, using 8, 16, 32 or all features. Its settings were fixed and not tuned: 300 trees
of depth 3, learning rate 0.05, and class weighting. The first experiments also used logistic
regression, a random forest and a support vector machine. XGBoost is a benchmark only; the
language model remains the model of interest.

\paragraph{Fusion.} To test whether the language model carries information that XGBoost does not
use, XGBoost is also trained on exactly the features of the best language-model setting, and the
two models' probabilities are combined in three ways: their mean, the mean of their ranks, and a
logistic regression on the two logits (stacking). The stacker is fitted on the validation split,
so calibration and test stay untouched and the conformal layer can be applied to the combined
score as to any other (\texttt{experiments/hybrid\_xgb\_llm.py}).

\section{Conformal triage layer}
\label{sec:methods-triage}

\subsection{Three-way decision rule}

Let $s(x)$ be a score where higher means more likely malignant. With two thresholds $t < u$:
\[
  \text{decision}(x) =
  \begin{cases}
    \text{auto-benign}    & s(x) \le t, \\
    \text{auto-malignant} & s(x) \ge u, \\
    \text{refer}          & t < s(x) < u.
  \end{cases}
\]
A referred nodule follows the usual clinical path (radiologist review or biopsy). The model is not
changed by this layer; only $t$ and $u$ are chosen, on the calibration split.

\subsection{Risks and Learn-then-Test calibration}

Two clinical risks are controlled:

\begin{itemize}
  \item $R_1$: the share of malignant nodules among the auto-benign ones is at most
        $\alpha_1 = 5\%$ (cancers wrongly cleared);
  \item $R_2$: the share of benign nodules among the auto-malignant ones is at most
        $\alpha_2 = 20\%$ (unneeded biopsies or treatment).
\end{itemize}

The thresholds are chosen with Learn then Test (LTT)~\citep{angelopoulos2021ltt}, which turns
threshold selection into a sequence of hypothesis tests. For a candidate $t$, let $n$ be the
number of calibration nodules with $s \le t$ and $k$ the number of malignant ones among them. The
null hypothesis ``the true risk exceeds $\alpha_1$'' is tested with the binomial p-value
\[
  p(t) = P\big(\mathrm{Bin}(n, \alpha_1) \le k\big).
\]
Candidates are fixed in advance and tested from the safest to the most permissive. Testing stops
at the first $p(t) > \delta$, and the last certified candidate is kept (fixed-sequence testing,
which needs no multiplicity correction, because the order does not depend on the calibration
data). The binomial test is valid because, given $n$, the calibration nodules with $s \le t$ are
an independent sample of the nodules below the threshold, so the number of cancers among them is
$\mathrm{Bin}(n, R_1(t))$. The guarantee is of the high-probability kind: with probability at least
$1 - \delta$ over the calibration draw, the true risk of the chosen threshold is at most
$\alpha_1$; the share observed on one test set can still exceed it. $u$ is chosen the same way with $\alpha_2$. With
$\delta = 0.05$ for each risk, both guarantees hold together with probability at least
$1 - 2\delta = 0.90$ over the draw of the calibration set, provided calibration and future nodules
are exchangeable (Section~\ref{sec:data-official}). If no candidate is certified, that side of the
rule makes no automatic decisions.

The test requires a group that is both large and clean. With $\alpha_1 = \delta = 0.05$, an
auto-benign group needs at least 59 nodules with no cancer, 93 nodules with one cancer and 153
nodules with three cancers to be certified.

For the softmax probability the candidates are $t \in \{0.05, 0.10, \dots, 0.45\}$ and
$u \in \{0.95, 0.90, \dots, 0.55\}$. Any order fixed without the calibration data keeps the
guarantee; as an exploratory variant, the candidates were also ordered by their $p$-values on the
validation split (Section~\ref{sec:results-ltt}).

\subsection{Uncertainty scores}
\label{sec:methods-scores}

\paragraph{Softmax probability.} The model's probability of malignancy for the unchanged prompt.

\paragraph{Perturbation votes.} Following the consistency-based uncertainty literature for
language models~\citep{ashby2026,noorani2025}, the original design of this work used an
agreement signal: the same nodule is presented to the unchanged model $m$ times in slightly
different ways, and the score is the number of queries $v$ answered malignant (probability
$\ge 0.5$). The first query is always the exact training prompt. The candidates for $m = 10$ are
$t \in \{0, \dots, 4\}$ and $u \in \{10, \dots, 6\}$. Two kinds of perturbations were used:

\begin{itemize}
  \item \textbf{Run 1 (mixed):} each query shuffled the feature order, dropped one feature with
        probability 0.5 and used one of three sentence templates.
  \item \textbf{Run 2 (format-preserving):} after run 1 showed that the model reacts to the format
        rather than to the content (Section~\ref{sec:results-ltt}), the template and order were
        kept. \emph{Jitter votes} add Gaussian noise to every standardised value (9 queries per
        noise level, standard deviation 0.05, 0.1 or 0.2, the level chosen on validation by a
        rule fixed in advance). \emph{Drop votes} remove one feature at a time (8 queries).
\end{itemize}

\subsection{Comparison with split conformal prediction}
\label{sec:methods-splitcp}

The same scores were also used with classical split conformal prediction~\citep{vovk2005}. The
nonconformity of a label is one minus its predicted probability, and a label enters the prediction
set when its nonconformity is below a threshold calibrated so that the true label is in the set
for at least $1-\alpha$ of the nodules. A set with one label counts as an automatic decision; a set
with both labels sends the nodule to a clinician. In the \emph{marginal} form one threshold serves
both labels; in the class-conditional (\emph{Mondrian}) form each label is calibrated on its own
nodules, so the guarantee holds for benign and malignant nodules separately. Unlike LTT, these
guarantees hold in expectation over the calibration and test draws, not with a stated
probability. On the official split
both forms were run at $\alpha = 0.10$. On the pooled split the marginal form was run at
$\alpha = 0.10$ and $0.05$, and the Mondrian form with $\alpha = 0.10$ for benign and $0.05$ for
malignant nodules: then at most about 5\% of cancers receive the single label ``benign'', a
guarantee on missed cancers like the miss-rate variant below. Neither form bounds the share of
cancers \emph{among} the single-label benign sets, which is what $R_1$ requires.

\subsection{Exploratory miss-rate variant}
\label{sec:methods-miss}

$R_1$ is a statement about each auto-benign decision. A weaker, programme-level alternative bounds
the \emph{miss rate}: the share of all cancers sent to auto-benign, as in conformal risk
control~\citep{angelopoulos2024crc}. The same binomial test is used with $n$ equal to the number of
calibration cancers, on the grid $t \in \{0.01, 0.02, \dots, 0.50\}$, at 5\% and 10\%. It is a valid
LTT procedure with a different risk: the number of calibration cancers below $t$ is
$\mathrm{Bin}(n, m(t))$, where $m(t)$ is the miss rate. It was added after the main results, so it
is reported as exploratory.

\section{Evaluation protocol}
\label{sec:methods-eval}

\begin{itemize}
  \item \textbf{Roles of the splits.} Training fits the model and all preprocessing; validation
        selects the epoch and the noise level; calibration is used only for the conformal
        thresholds; test only for reporting, with the rules locked.
  \item \textbf{Classification.} The main measure is the area under the ROC curve (AUC), computed
        from the predicted probability. Accuracy and sensitivity use the 0.5 threshold.
  \item \textbf{Comparing two runs.} Differences in AUC are given with a 95\% paired bootstrap
        interval: the same nodules are resampled 2,000 times and the AUC difference of the two
        runs is recomputed on each resample.
  \item \textbf{Separating errors.} How well a confidence score separates correct from incorrect
        predictions is measured by the AUC of that score for the event ``prediction correct''.
  \item \textbf{Exchangeability.} The classifier two-sample test of Section~\ref{sec:data-official}.
  \item \textbf{Training noise.} The main comparison (with and without the TI-RADS-motivated
        features) is repeated with three seeds (42, 1, 2), matched across the two settings, and
        reported as mean $\pm$ standard deviation, per-seed paired differences, and the interval
        of the seed-averaged probabilities. The bootstrap intervals alone do not cover this
        variation, because they hold the trained model fixed.
  \item \textbf{Conformal coverage over re-splits.} A single test split shows one realisation of
        a guarantee that holds on average. With the model fixed, calibration and test are pooled
        and randomly re-split 4,000 times (same sizes and class proportions), and the conformal
        thresholds are recomputed each time (\texttt{experiments/conformal\_checks.py}).
\end{itemize}

Unless stated otherwise, a run uses seed 42 and is a single training run.

\chapter{Experiments and results}
\label{ch:results}

The experiments were run in the order of Table~\ref{tab:results-timeline}, and each answered a
question raised by the one before. This chapter is organised by question instead. Most settings
were trained once, with seed 42. Training noise turned out to be large
(Section~\ref{sec:results-seeds}), so the main comparison, with and without the TI-RADS-motivated
features, was repeated with three seeds, and single-run differences below about 0.03--0.05 AUC are
not interpreted.

\section{Timeline of experiments}
\label{sec:results-timeline}

\begin{table}[ht]
  \centering
  \small
  \caption{Experiments. Records are folders in \texttt{runs/}; the consolidated report is
           \texttt{reports/consolidated\_report\_2026-10-12.md}.}
  \label{tab:results-timeline}
  \begin{tabular}{lp{4.3cm}p{6.2cm}}
    \toprule
    date & experiment & outcome \\
    \midrule
    2 Oct  & Triage layer, run 1: mixed perturbations (official split, v1, 8 features, Qwen2.5-1.5B) & nothing certified; perturbations push the model off its training format \\
    2 Oct  & Triage layer, run 2: format-preserving perturbations & nothing certified; votes no better than softmax; calibration and official test not exchangeable \\
    7 Oct  & Error analysis and miss-rate variant (no new training) & errors are stable across queries; miss rate certifiable at 5\% for few nodules \\
    8 Oct  & Backbone sweep on the official split (Qwen2.5-1.5B retrained, Qwen2.5-7B) & 7B not better \\
    8 Oct  & Pooled split; image quality control; XGBoost & exchangeability restored; feature problems found \\
    8--9 Oct & v2 features; LLM runs A--F on the pooled split & v2 helps (+0.05 AUC); format and size results \\
    9 Oct  & Llama-2-7B with the settings of Ra et al.\ (run G); rich format (run H) & neither helps \\
    9--10 Oct & TI-RADS-motivated features (runs I, J, K) & largest gain of all changes; seventh feature and rich format add nothing \\
    11 Oct & Qwen2.5-7B at a lower learning rate, with TI-RADS features (run L) & clearly worse than 1.5B \\
    12 Oct & Seed repeats with and without TI-RADS features (runs M--P); fusion with XGBoost; conformal checks & TI-RADS gain confirmed; the language model adds nothing to XGBoost; Mondrian coverage holds over re-splits \\
    \bottomrule
  \end{tabular}
\end{table}

\section{Learning from radiomics prompts}
\label{sec:results-llm}

Yes, to a moderate degree. On the official split, with eight v1 features, Qwen2.5-1.5B reached
test AUC 0.672 (validation 0.704, calibration 0.707), against 0.648--0.663 for a logistic
regression, a random forest and a support vector machine. A retrained copy of the same model gave
0.667 on test, and its scores agreed closely with the first model (Spearman correlation 0.97).
Validation and calibration AUCs above the test AUC are consistent with the shift described in
Section~\ref{sec:data-official}.

Table~\ref{tab:results-llm} lists the runs on the pooled split. With the radiomics features alone
the best settings reach a test AUC of 0.75--0.76; with the TI-RADS-motivated features, 0.785 on
average over three seeds.

\begin{table}[ht]
  \centering
  \small
  \setlength{\tabcolsep}{4pt}
  \caption{Language-model runs on the pooled split (16 radiomics features, 5 epochs; learning rate
           $2 \times 10^{-4}$, LoRA on Q/V and seed 42 unless stated). Sensitivity at threshold
           0.5 on test. (f): near-duplicate filter; +T6, +T7: six or seven TI-RADS-motivated
           features. Run G uses the settings of Ra et al.\ (learning rate $10^{-4}$, LoRA on all
           linear layers); run L uses learning rate $5 \times 10^{-5}$.}
  \label{tab:results-llm}
  \begin{tabular}{llllrrrrr}
    \toprule
    run & model & features & format & val & cal & test & sens. & minutes \\
    \midrule
    A & Qwen2.5-1.5B & v2            & decimal    & 0.753 & 0.757 & 0.750 & 0.66 & 26 \\
    B & Qwen2.5-1.5B & v1 (f)        & decimal    & 0.709 & 0.705 & 0.699 & 0.67 & 27 \\
    C & Qwen2.5-7B   & v2 (f)        & decimal    & 0.708 & 0.705 & 0.713 & 0.00 & 90 \\
    D & Qwen2.5-1.5B & v2 (f)        & decimal    & 0.750 & 0.752 & 0.754 & 0.71 & 26 \\
    E & Qwen2.5-1.5B & v2 (f)        & z-score    & 0.729 & 0.756 & 0.759 & 0.65 & 37 \\
    F & Qwen2.5-1.5B & v2 (f)        & percentile & 0.709 & 0.713 & 0.734 & 0.65 & 37 \\
    G & Llama-2-7B   & v2 (f)        & decimal    & 0.705 & 0.709 & 0.731 & 0.66 & 182 \\
    H & Qwen2.5-1.5B & v2 (f)        & rich       & 0.732 & 0.711 & 0.744 & 0.64 &  83 \\
    I & Qwen2.5-1.5B & v2 (f) +T6    & decimal    & 0.795 & 0.812 & 0.804 & 0.65 & 36 \\
    J & Qwen2.5-1.5B & v2 (f) +T7    & decimal    & 0.786 & 0.802 & 0.781 & 0.39 & 37 \\
    K & Qwen2.5-1.5B & v2 (f) +T6    & rich       & 0.803 & 0.785 & 0.801 & 0.64 & 132 \\
    L & Qwen2.5-7B   & v2 (f) +T6    & decimal    & 0.751 & 0.712 & 0.693 & 0.56 & 123 \\
    \midrule
    M & I, seed 1    & v2 (f) +T6    & decimal    & 0.771 & 0.777 & 0.768 & 0.66 & 37 \\
    N & I, seed 2    & v2 (f) +T6    & decimal    & 0.787 & 0.794 & 0.783 & 0.69 & 37 \\
    O & D, seed 1    & v2 (f)        & decimal    & 0.706 & 0.695 & 0.712 & 0.69 & 36 \\
    P & D, seed 2    & v2 (f)        & decimal    & 0.723 & 0.715 & 0.713 & 0.65 & 36 \\
    \bottomrule
  \end{tabular}
\end{table}

Training was unstable in several runs. In the early epochs validation sensitivity at 0.5 was
sometimes close to 0 (all nodules called benign) and recovered later. Selecting the epoch by
validation AUC hides this from the final AUC, but not from threshold-based measures such as
sensitivity (run J ends at 0.39).

\section{How large is the training noise?}
\label{sec:results-seeds}

Larger than the effects of most changes. The best setting (run I) was trained again with seeds 1
and 2 (runs M and N), and the reference without TI-RADS features (run D) likewise (runs O and P).
Only the initialisation of the adapters and the order of the training batches change; the split
and the selected features stay the same.

\begin{itemize}
  \item The three seeds of the best setting give test AUC 0.804, 0.768 and 0.783: mean 0.785,
        standard deviation 0.018. Seed 42 was the upper end. Runs M and I differ significantly on
        every split (test interval $[-0.058, -0.014]$).
  \item The three seeds without TI-RADS features give 0.754, 0.712 and 0.713; again seed 42 was
        the best.
  \item A repeat of run H with the same seed, where only the padding length differed, moved by
        0.046 on test.
\end{itemize}

A paired bootstrap interval holds the trained model fixed and resamples nodules, so it does not
contain this variation. Comparisons between single runs therefore understate the uncertainty, and
the result that matters most for the thesis, the effect of the TI-RADS-motivated features, is
reported over matched seeds (Section~\ref{sec:results-clinical}).

\section{Do better features help?}
\label{sec:results-features}

Yes, clearly. Runs B and D differ only in the feature set: the test AUC rises from 0.699 with v1
to 0.754 with v2. The paired bootstrap interval for run A against run B is $[+0.013, +0.089]$ on
test and $[+0.015, +0.090]$ on calibration, so the gain of about 0.05 is not a chance
fluctuation of the test set.

XGBoost shows the same picture (Table~\ref{tab:results-xgb}). With v1, adding features helps
slowly; with v2, 16 features already give 0.78--0.81.

\begin{table}[ht]
  \centering
  \caption{XGBoost test AUC on the pooled split, by number of mRMR features. ``Filtered'':
           near-duplicates ($|r| > 0.95$) removed first; the last column gives the number of
           features left.}
  \label{tab:results-xgb}
  \begin{tabular}{lrrrr}
    \toprule
    features & 8 & 16 & 32 & all \\
    \midrule
    v1           & 0.695 & 0.723 & 0.738 & 0.793 (67) \\
    v1, filtered & 0.691 & 0.717 & 0.743 & 0.797 (44) \\
    v2           & 0.723 & 0.780 & 0.803 & 0.800 (102) \\
    v2, filtered & 0.723 & \textbf{0.805} & 0.796 & 0.800 (54) \\
    \bottomrule
  \end{tabular}
\end{table}

The near-duplicate filter lifts XGBoost with 16 v2 features from 0.780 to 0.805, because the
16 slots are no longer spent on copies. For the language model it makes no difference (runs A and
D, within 0.005).

On the official split, more features had seemed useless: classical models gained 0.05--0.10 AUC
on validation by going from 8 to 67 features but not on the official test set. On the pooled split
the gain carries over to the test set, so the earlier result was a consequence of the shift
between the official sources, not of the features.

\section{Does the number format matter?}
\label{sec:results-format}

Saying what the number is does not matter; how extreme the number is does; and adding what the
value means for the nodule does not help.

\begin{table}[ht]
  \centering
  \caption{AUC differences between formats, 95\% paired bootstrap intervals.}
  \label{tab:results-format}
  \begin{tabular}{llrr}
    \toprule
    comparison & what changes & calibration & test \\
    \midrule
    E $-$ D & z-score wording               & $[-0.016, +0.022]$ & $[-0.017, +0.029]$ \\
    F $-$ D & percentile values             & $[-0.075, -0.005]$ & $[-0.057, +0.022]$ \\
    H $-$ D & rich format                   & $[-0.074, -0.011]$ & $[-0.041, +0.023]$ \\
    K $-$ I & rich format, with TI-RADS     & $[-0.049, -0.005]$ & $[-0.031, +0.022]$ \\
    \bottomrule
  \end{tabular}
\end{table}

\begin{itemize}
  \item The z-score wording (run E) and the bare decimal (run D) perform the same.
  \item Percentiles (run F) lose about 0.04 AUC on validation and calibration; on test the
        interval includes zero. A percentile maps a z-score of 2 and a z-score of 5 to nearly the
        same value (98 and 100), so the model loses the information that a value is extreme.
        The loss suggests that the model uses the magnitude of the numbers, not only their order.
  \item The \emph{rich} format (run H) keeps the z-score and adds the percentile, a band
        (``lower than typical'') and what the value means (``meaning a less round, more irregular
        outline''). It does not help: test AUC 0.744 against 0.754, and 0.04 lower on calibration,
        where the interval excludes zero. A repeat with the same seed gave 0.698.
  \item With the TI-RADS-motivated features (run K against run I), the rich format is equal on
        validation and test and slightly worse on calibration. Its prompts are up to three times
        as long, and training took 132 minutes instead of 36. Training was stable, so the weak
        rich runs without these features were not caused by prompt length alone.
\end{itemize}

Plain-language meaning gives the model nothing it does not already get from the numbers. The
decimal format is kept.

\section{Does a larger backbone help?}
\label{sec:results-size}

No: neither a larger Qwen model, at two learning rates, nor the backbone and settings of Ra et al.\
beat Qwen2.5-1.5B.

\begin{itemize}
  \item On the official split, Qwen2.5-7B matched the 1.5B model: test AUC 0.669 with 8 features
        and 0.666 with 16, against 0.667 for 1.5B.
  \item On the pooled split, Qwen2.5-7B (run C) was worse than 1.5B (run D): 0.713 against 0.754,
        paired interval $[-0.069, -0.004]$ on test.
  \item With the backbone and settings of Ra et al.\ (run G: Llama-2-7B, learning rate
        $10^{-4}$, LoRA on all linear layers), the test AUC is 0.731, again below run D. The
        paired intervals against run D are $[-0.070, -0.016]$ on calibration and
        $[-0.052, +0.007]$ on test; against Qwen2.5-7B (run C) they include zero.
  \item With the TI-RADS-motivated features and a four times lower learning rate (run L),
        Qwen2.5-7B reaches 0.693, far below the 1.5B model with the same features (run I, 0.804;
        intervals exclude zero on every split).
\end{itemize}

The first 7B training did not settle. Across the five epochs, validation sensitivity at threshold
0.5 went 0.00, 0.01, 0.93, 0.04, 0.62; gradient norms reached 200; the best epoch by validation AUC
was epoch 2, and the selected model calls every test nodule benign at 0.5 although its ranking
still has AUC 0.713. The learning rate of $2 \times 10^{-4}$, chosen for the 1.5B model, was the
obvious suspect, so run L used $5 \times 10^{-5}$. This calmed the swings (sensitivity 0.85, 0.30,
0.31, 0.46, 0.58) but not the result: gradient norms still reached 100--200, validation AUC crept
up to 0.753 while the learning rate decayed to zero, and calibration and test fell well below
validation (0.712 and 0.693). The TI-RADS features, which lift the 1.5B model by about 0.06, gave
Qwen2.5-7B nothing (run L against run C: test interval $[-0.062, +0.023]$). Run G took 182
minutes on one L4 GPU, seven times as long as Qwen2.5-1.5B, for a lower AUC.

Within this setup (QLoRA on Q and V, about 2,000 training prompts, 5 epochs), a larger backbone
does not help, and Qwen2.5-1.5B stays the main model.

\section{Do TI-RADS-motivated features help?}
\label{sec:results-clinical}

Yes, for both models, and more than any other change tried. Table~\ref{tab:results-clinical-uni}
shows each feature on its own, on the training split.

\begin{table}[ht]
  \centering
  \small
  \caption{TI-RADS-motivated features on their own: AUC on the training split (above 0.5: higher
           in malignant nodules) and the direction TI-RADS leads one to expect.}
  \label{tab:results-clinical-uni}
  \begin{tabular}{llrl}
    \toprule
    category & feature & train AUC & as expected? \\
    \midrule
    shape          & taller-than-wide ratio & 0.658 & yes (higher in cancers) \\
    margin         & solidity               & 0.367 & yes (lower in cancers) \\
    composition    & anechoic fraction      & 0.409 & yes (benign nodules more cystic) \\
    echogenicity   & echogenicity ratio     & 0.552 & no (expected below 0.5) \\
    margin         & margin sharpness       & 0.563 & no (expected below 0.5) \\
    echogenic foci & punctate foci density  & 0.487 & no signal \\
    \bottomrule
  \end{tabular}
\end{table}

Three features behave as TI-RADS predicts, and the taller-than-wide ratio alone is as
discriminative as the best single radiomics features. The echogenicity ratio points the wrong way:
it correlates $-0.78$ with the anechoic fraction, so it mostly measures fluid. Cystic nodules,
mostly benign, are dark because of their fluid, while TI-RADS judges echogenicity on the solid
part. A seventh feature computed on the non-fluid pixels only (run J) removed the reversed
direction but added no signal; with it the language model was slightly worse (test interval
$[-0.043, -0.001]$ against run I), so the six-feature set is kept. Margin sharpness may measure
contrast rather than blur, and the simple spot detector finds no signal for microcalcifications in
speckled images.

Added to the 16 radiomics features, the six features raise XGBoost from 0.805 to 0.831 on test
(Table~\ref{tab:results-clinical-xgb}), with the same gain on validation and calibration. Adding the
remaining 38 radiomics features instead gives nothing (0.796), so the gain comes from information
that the radiomics features do not contain. The six features alone reach 0.775.

\begin{table}[ht]
  \centering
  \small
  \caption{XGBoost with and without the TI-RADS-motivated features (pooled split, radiomics
           features after the near-duplicate filter).}
  \label{tab:results-clinical-xgb}
  \begin{tabular}{lrrrrr}
    \toprule
    features & number & val & cal & test & \multicolumn{1}{p{2.4cm}}{\raggedleft cancers in lowest 15\% of cal} \\
    \midrule
    radiomics, mRMR 16            & 16 & 0.799 & 0.823 & 0.805 & 4/78 \\
    radiomics, mRMR 16 + TI-RADS  & 22 & 0.828 & 0.848 & 0.831 & 2/78 \\
    TI-RADS only                  &  6 & 0.786 & 0.787 & 0.775 & 6/78 \\
    radiomics, all                & 54 & 0.805 & 0.816 & 0.796 & 2/78 \\
    radiomics, all + TI-RADS      & 60 & 0.824 & 0.836 & 0.826 & 3/78 \\
    \bottomrule
  \end{tabular}
\end{table}

The language model gains more. Because a single run can be off by 0.03--0.05
(Section~\ref{sec:results-seeds}), the comparison was repeated with three matched seeds: each seed
trains the model once with and once without the six features
(Table~\ref{tab:results-clinical-seeds}).

\begin{table}[ht]
  \centering
  \small
  \setlength{\tabcolsep}{4pt}
  \caption{Language model with and without the TI-RADS-motivated features, three matched seeds
           (42, 1, 2). Mean $\pm$ standard deviation over seeds; the last row averages the three
           seeds' probabilities per nodule and gives a 95\% paired bootstrap interval.}
  \label{tab:results-clinical-seeds}
  \begin{tabular}{lrrr}
    \toprule
    & val & cal & test \\
    \midrule
    with TI-RADS (I, M, N)        & $0.785 \pm 0.012$ & $0.794 \pm 0.017$ & $0.785 \pm 0.018$ \\
    without (D, O, P)             & $0.727 \pm 0.022$ & $0.721 \pm 0.029$ & $0.726 \pm 0.024$ \\
    difference per seed, mean (min) & $+0.058$ ($+0.045$) & $+0.074$ ($+0.059$) & $+0.059$ ($+0.050$) \\
    seed-averaged, 95\% interval  & $[+0.018, +0.092]$ & $[+0.033, +0.100]$ & $[+0.019, +0.094]$ \\
    \bottomrule
  \end{tabular}
\end{table}

The gain is positive for every seed and every split, at least $+0.045$, and the test interval of
each matched pair excludes zero (seed 42: $[+0.015, +0.084]$; seed 1: $[+0.013, +0.098]$; seed 2:
$[+0.026, +0.112]$). Seed 42 gave the best run on both sides, so the single-run comparison of
runs I and D (+0.050) did not overstate the gain: on average it is about $+0.06$, slightly more
than for XGBoost. Training with the features was also steadier: validation AUC rose over the
epochs in all three seeds.

\subsection{Peritumoral radiomics}
\label{sec:results-peri}

Since information about the surroundings helped, the PyRadiomics features were also computed in a
ring of tissue around the nodule (width half of the nodule's equivalent radius, at least 5
pixels; all families except shape; 93 features, 49 after the near-duplicate filter;
\texttt{conformal\_triage/peritumoral.py}). The rule was fixed before running: the ring features
would be given to the language model only if adding eight of them (mRMR) to XGBoost improved the
calibration AUC with an interval above zero.

\begin{table}[ht]
  \centering
  \small
  \caption{XGBoost with peritumoral (ring) features, pooled split.}
  \label{tab:results-peri}
  \begin{tabular}{lrrrr}
    \toprule
    features & number & val & cal & test \\
    \midrule
    nodule mRMR 16 + TI-RADS                    & 22 & 0.828 & 0.848 & 0.831 \\
    \quad + ring mRMR 8                         & 30 & 0.832 & 0.851 & 0.834 \\
    \quad + all ring features                   & 71 & 0.831 & 0.853 & 0.832 \\
    joint mRMR 16 (nodule and ring) + TI-RADS   & 22 & 0.836 & 0.840 & 0.821 \\
    ring mRMR 16 only                           & 16 & 0.699 & 0.671 & 0.693 \\
    \bottomrule
  \end{tabular}
\end{table}

Adding the ring features changes the calibration AUC by $+0.002$, with a 95\% interval from
$-0.009$ to $+0.014$, and the test AUC by $+0.003$ (from $-0.010$ to $+0.016$). By the rule they
were not passed to the language model. The ring carries signal on its own (test 0.693), but it is redundant with the nodule and
TI-RADS-motivated features: the ring features mRMR selects are mostly brightness statistics, close
to what the echogenicity ratio already compares. What helped was not more features about the
nodule's neighbourhood but the specific information TI-RADS points to, such as orientation and
outline, that standard radiomics does not measure.

\section{How does the language model compare with XGBoost?}
\label{sec:results-xgb}

XGBoost is better with the same features, by about 0.05 AUC, and the language model adds nothing
to it.

\begin{itemize}
  \item Official split, 8 v1 features: XGBoost 0.696, language model 0.667; with 16 features
        XGBoost reaches 0.709 while the 7B language model stays at 0.666.
  \item Pooled split, 16 v2 features after filtering: XGBoost 0.805, language model 0.754
        (run D).
  \item Pooled split, the same 22 features as run I: XGBoost 0.835, language model
        $0.785 \pm 0.018$ over three seeds.
\end{itemize}

Run I alone (0.804) had suggested that the gap shrinks to 0.03 with the TI-RADS features; over
seeds it stays near 0.05.

\subsection{Fusion of the two models}
\label{sec:results-fusion}

A smaller AUC does not rule out that the language model carries information XGBoost misses. To
test this, XGBoost was trained on exactly the 22 features of run I with the same preprocessing,
and the two models' probabilities were combined (\texttt{experiments/hybrid\_xgb\_llm.py}). The
language model is the average of its three seeds.

\begin{table}[ht]
  \centering
  \caption{XGBoost, the language model and their combinations, on the 22 features of run I.
           Stacking is a logistic regression on the two logits, fitted on validation only.}
  \label{tab:results-fusion}
  \begin{tabular}{lrrr}
    \toprule
    model & val & cal & test \\
    \midrule
    XGBoost                                  & 0.827 & 0.848 & \textbf{0.835} \\
    language model, average of 3 seeds       & 0.792 & 0.803 & 0.793 \\
    mean of the two probabilities            & 0.820 & 0.838 & 0.824 \\
    mean of the two ranks                    & 0.819 & 0.837 & 0.824 \\
    stacking                                 & 0.827 & 0.848 & 0.834 \\
    \bottomrule
  \end{tabular}
\end{table}

\begin{itemize}
  \item The language model is below XGBoost on every split (test difference $-0.042$, interval
        $[-0.067, -0.017]$).
  \item Averaging the two is slightly worse than XGBoost alone (test $-0.011$, interval
        $[-0.023, +0.001]$).
  \item The stacker gives XGBoost a weight of 0.99 and the language model $-0.04$: it learns to
        ignore the language model, and its AUC equals that of XGBoost.
  \item The two models' probabilities correlate 0.87--0.88 on every split. Two seeds of the
        language model correlate 0.91, so the language model differs from XGBoost hardly more
        than from itself.
\end{itemize}

On these features the fine-tuned language model reproduces a weaker version of the tabular
model's ranking and carries no information that XGBoost does not already use.

\section{Can Learn-then-Test certify automatic decisions?}
\label{sec:results-ltt}

No. In none of the runs, with any score, did the planned LTT procedure certify a single automatic
decision at $\alpha_1 = 5\%$ or $\alpha_2 = 20\%$, and neither did it for XGBoost. This section
explains why, starting with the perturbation votes of the original design.

\subsection{Mixed perturbations push the model off its format}

In run 1 (official split, validation nodules), shuffling the feature order alone changed 29\% of
the answers. Paraphrased templates changed 47--48\% of the answers and were \emph{never}
classified malignant, in any split. No nodule received more than 6 of 10 malignant votes, so the
auto-malignant side was impossible by construction. The model also depended on feature position:
shuffling changed only 10\% of the answers when the first feature stayed first, against 24--39\%
otherwise. These perturbations measured the distance from the training format, not the
uncertainty about the nodule.

\subsection{Format-preserving votes are not better than softmax}

With jitter and drop-one perturbations (run 2), the model stays in distribution: noise and
dropping change 5--12\% of the answers and the AUC of the perturbed queries stays within 0.02 of
the base prompt. But the votes separate correct from incorrect predictions no better than the
softmax probability (Table~\ref{tab:results-errsep}). Jitter votes are worse on calibration and
test; drop votes are level with softmax.

\begin{table}[ht]
  \centering
  \caption{AUC for separating correct from incorrect predictions (run 2, official split).}
  \label{tab:results-errsep}
  \begin{tabular}{lrrr}
    \toprule
    confidence score & validation & calibration & test \\
    \midrule
    softmax confidence          & 0.602 & 0.629 & 0.577 \\
    jitter vote agreement (SD 0.2) & 0.560 & 0.579 & 0.533 \\
    drop vote agreement         & 0.620 & 0.652 & 0.571 \\
    \bottomrule
  \end{tabular}
\end{table}

\subsection{The model is consistently wrong}

The design assumed that the model would be unstable on the nodules it gets wrong, so that their
votes would split and they would be referred. The opposite happened. On calibration, 19.4\% of the
nodules with 0 of 10 malignant votes are malignant, and 45.7\% of those with 10 of 10 are benign.
Between 56\% and 69\% of the base prompt's errors are unanimous: all ten queries give the same
wrong answer. Consistency measures how stable an answer is, not whether it is correct, and no vote
threshold can separate these errors.

\subsection{Nothing certified, for any model}

Table~\ref{tab:results-ltt} shows the cleanest auto-benign group available in each setting. All
are far from what LTT needs: 59 nodules without a cancer, 93 with one or 124 with two
(Section~\ref{sec:methods-triage}).

\begin{table}[ht]
  \centering
  \small
  \caption{Cleanest auto-benign groups on calibration (524 nodules, 182 cancers on the pooled
           split).}
  \label{tab:results-ltt}
  \begin{tabular}{p{4.2cm}p{7.4cm}}
    \toprule
    setting & best auto-benign group \\
    \midrule
    official, jitter votes  & $v = 0$: 165 nodules, 19.4\% malignant \\
    official, drop votes    & $v = 0$: 175 nodules, 20.0\% malignant \\
    official, softmax       & not certified at any candidate \\
    pooled, run E           & lowest 15\%: 78 nodules, 1 malignant; $P \le 0.20$: 90 nodules, 3 malignant \\
    pooled, run I           & lowest 15\%: 78 nodules, 8 malignant; $P \le 0.05$: 4 nodules, none malignant \\
    pooled, LLM, mean of 3 seeds & lowest 10\%, 15\%, 20\%: 4/52, 7/78, 10/104 malignant \\
    pooled, XGBoost (22 features) & lowest 10\%, 15\%, 20\%: 2/52, 2/78, 3/104 malignant \\
    \bottomrule
  \end{tabular}
\end{table}

Run E comes closest among the language models: its 78 lowest-scored calibration nodules contain
one cancer (1.3\%), well below the 5\% target, but with one cancer 93 nodules are needed. The best
setting is much dirtier there: 8 cancers among the lowest 78 for run I and 7 for the average of its
three seeds. XGBoost, with the TI-RADS
features, has the cleanest tail of all models (2 cancers among the lowest 78, 3 among the lowest
104), and still nothing is certified: two cancers need 124 nodules, three need 153.

The auto-malignant side shows a second obstacle, which lies in the procedure rather than the
model. For XGBoost, the candidate $u = 0.85$ would pass on its own (54 calibration nodules, 4
benign, $p = 0.010$ against $\alpha_2 = 20\%$). The planned fixed sequence, however, starts at the
strictest threshold, $u = 0.95$, where only 7 nodules fall and 3 of them are benign; it fails
there, stops, and never reaches 0.85. Run I has the same pattern ($u = 0.85$: 50 nodules, 4
benign, $p = 0.018$; $u = 0.95$: 13 nodules, $p = 0.23$). The strictest candidates hold too few
nodules to ever pass, so a sequence that starts with them stops at once.

Ordering the candidates by their $p$-values on the validation split, which is independent of
calibration and therefore keeps the guarantee, avoids starting at a near-empty threshold. This
certifies auto-malignant at $u = 0.85$ for run I (test: 43 nodules, 9.8\%, of which 5 benign,
11.6\%) and at $u = 0.90$ for the language model averaged over three seeds (test: 28 nodules,
6.4\%, none benign). It certifies nothing for runs D and E, nothing for XGBoost (whose validation
order still puts 0.95 first), and nothing on the auto-benign side. Because the ordering rule was
chosen after run I was seen, this result is exploratory: it has to be fixed in advance and
confirmed on another split before it can be claimed.

\subsection{A design fixed in advance, over re-splits}
\label{sec:results-ltt-design}

The two obstacles suggest a design that keeps the guarantee and can pass. It was written down and
committed before it was run; the code is in
\texttt{experiments/\allowbreak ltt\_design\_resplits.py}.

\begin{itemize}
  \item \textbf{Candidates} at quantiles of the validation scores, so that each holds a known
        share of the nodules: auto-benign below the 5\%, 10\%, \dots, 30\% quantile, auto-malignant
        above the 95\%, 90\%, \dots, 70\% quantile.
  \item \textbf{Order} by the candidates' binomial $p$-values on validation, smallest first.
  \item \textbf{Test} on calibration with the fixed sequence, $\alpha_1 = 5\%$,
        $\alpha_2 = 20\%$, $\delta = 0.05$, as planned.
\end{itemize}

Keeping the models fixed, calibration and test were re-split 2,000 times as in
Section~\ref{sec:results-splitcp}, and the new design was compared with the planned one
(Table~\ref{tab:results-ltt-design}).

\begin{table}[ht]
  \centering
  \small
  \setlength{\tabcolsep}{4pt}
  \caption{Share of 2,000 re-splits in which a side is certified, and, when the quantile design
           certifies auto-malignant, the share of test nodules decided and the share of benign
           nodules among them (on test, and on all 962 nodules above the threshold).}
  \label{tab:results-ltt-design}
  \begin{tabular}{lrrrrr}
    \toprule
    & \multicolumn{1}{p{1.5cm}}{\raggedleft auto-benign, either design}
    & \multicolumn{1}{p{1.6cm}}{\raggedleft auto-malig., planned}
    & \multicolumn{1}{p{1.6cm}}{\raggedleft auto-malig., quantile}
    & \multicolumn{1}{p{1.4cm}}{\raggedleft decided, test}
    & \multicolumn{1}{p{1.9cm}}{\raggedleft benign among them (test / all)} \\
    \midrule
    XGBoost (22 features) & 0\% & 0\%  & 38\% & 10.7\% & 14.9\% / 11.2\% \\
    LLM, mean of 3 seeds  & 0\% & 1\%  & \textbf{99\%} & 7.9\% & 7.4\% / 6.4\% \\
    LLM run I             & 0\% & 13\% & 22\% & 11.5\% & 18.4\% / 13.6\% \\
    LLM run D             & 0\% & 0\%  & 8\%  & 6.8\%  & 22.6\% / 13.8\% \\
    \bottomrule
  \end{tabular}
\end{table}

\begin{itemize}
  \item \textbf{Auto-benign at 5\% is never certified}, for any model or design. This side is out
        of reach of the current models; no choice of procedure changes that.
  \item \textbf{Auto-malignant at 20\% becomes attainable.} For the language model averaged over
        its three seeds, the new design certifies in 99\% of re-splits, against 1\% for the
        planned design. About 8\% of the nodules are then called malignant automatically, 6--7\%
        of them benign. On our split the rule is $u = 0.833$: 9.8\% of the test nodules, 14\% of
        them benign.
  \item \textbf{The order only helps where validation and calibration agree.} For XGBoost the
        validation order puts its top 5\% first, which on calibration contains 3 benign nodules
        among 7, so the sequence often stops there (38\%). The seed-averaged language model, whose
        scores are smoother, benefits most.
  \item \textbf{The bound holds.} The share of benign nodules among all 962 nodules above the
        chosen threshold never exceeds 20\% in any re-split. The test-only share can (run D):
        in a fixed pool, a re-split with a clean calibration part leaves a dirtier test part, so
        conditioning on certification biases the test share upwards.
\end{itemize}

Two caveats limit this result. The design was written after the exploratory result on run I, and
the models were selected on these data; the re-splits vary only which nodules calibrate and which
test. Confirmation needs new data: retraining on new splits, or a second dataset.

\subsection{A guarantee on missed cancers}

Bounding the share of \emph{all} cancers sent to auto-benign (Section~\ref{sec:methods-miss}) is
easier, because the evidence comes from all 182 calibration cancers. At 5\% it is certified in
every pooled run (Table~\ref{tab:results-miss}).

\begin{table}[ht]
  \centering
  \small
  \caption{Miss-rate guarantee at 5\%, pooled split. LLM avg.: the language model averaged over
           runs I, M and N; XGB: XGBoost on the 22 features of run I.}
  \label{tab:results-miss}
  \begin{tabular}{lrrrr}
    \toprule
    run & threshold & auto-benign, cal & auto-benign, test & cancers missed, test \\
    \midrule
    A & 0.16 & 14.3\% & 14.4\% & 5.3\% \\
    B & 0.22 &  4.0\% &  5.7\% & 2.6\% \\
    C & 0.09 & 10.3\% & 11.4\% & 2.6\% \\
    D & 0.16 & 16.8\% & 18.0\% & 8.6\% \\
    E & 0.21 & 19.1\% & 19.9\% & 5.9\% \\
    F & 0.26 &  9.2\% & 11.0\% & 3.3\% \\
    G & 0.17 &  6.9\% &  7.8\% & 0.7\% \\
    H & 0.20 &  4.2\% &  3.7\% & 0.7\% \\
    I & 0.09 &  6.1\% &  5.5\% & 0.7\% \\
    M & 0.09 &  5.2\% &  5.3\% & 1.3\% \\
    N & 0.12 &  9.5\% &  8.4\% & 2.6\% \\
    \midrule
    LLM avg. & 0.14 & 11.3\% & 12.1\% & 4.6\% \\
    XGB & 0.14 & 19.8\% & 23.7\% & 4.6\% \\
    \bottomrule
  \end{tabular}
\end{table}

The realised test miss rate ranges from 0.7\% to 8.6\%. A guarantee with $\delta = 0.05$ is
allowed to fail on some test sets, and the test set has only 152 cancers, so one cancer is 0.7
percentage points; one run above 5\% is not a contradiction. On the official split the same rule
behaved worse under the shift: at 10\%, 15.3\% of the official-test cancers were missed. XGBoost
sends twice as many nodules to auto-benign as the averaged language model at the same miss rate,
and the cancers among them are 7 of 104 (6.7\%) against 7 of 53 (13.2\%).

\subsection{Split conformal prediction on the pooled split}
\label{sec:results-splitcp}

Classical split conformal prediction (Section~\ref{sec:methods-splitcp}) gives what it promises on
the pooled split: with marginal $\alpha = 0.10$ the test coverage is 0.88--0.93, and with
$\alpha = 0.05$ it is 0.93--0.97. Marginal coverage, however, is carried by the benign majority:
at $\alpha = 0.10$ only 66--92\% of the test cancers have ``malignant'' in their set. The Mondrian
form protects the cancers directly (Table~\ref{tab:results-splitcp}).

\begin{table}[ht]
  \centering
  \small
  \caption{Mondrian split conformal prediction on the pooled test set ($\alpha = 0.10$ benign,
           $0.05$ malignant; thresholds from calibration). Auto-benign and auto-malignant are
           the single-label sets; the rest is referred. XGB: XGBoost on the 22 features of run I.}
  \label{tab:results-splitcp}
  \begin{tabular}{lrrrrrr}
    \toprule
    run & \multicolumn{1}{p{1.6cm}}{\raggedleft cancers covered}
        & \multicolumn{1}{p{1.3cm}}{\raggedleft auto-benign}
        & \multicolumn{1}{p{1.7cm}}{\raggedleft cancers in auto-benign}
        & \multicolumn{1}{p{1.5cm}}{\raggedleft auto-malignant}
        & \multicolumn{1}{p{1.7cm}}{\raggedleft benign in auto-malignant}
        & referred \\
    \midrule
    A & 0.901 & 22.8\% & 15.0\% & 19.6\% & 30.2\% & 57.5\% \\
    B & 0.914 & 17.1\% & 17.3\% & 17.8\% & 34.6\% & 65.1\% \\
    C & 0.941 & 18.7\% & 11.0\% & 16.4\% & 41.7\% & 64.8\% \\
    D & 0.908 & 21.7\% & 14.7\% & 17.4\% & 30.3\% & 61.0\% \\
    E & 0.934 & 22.8\% & 10.0\% & 16.7\% & 27.4\% & 60.5\% \\
    F & 0.928 & 18.5\% & 13.6\% & 15.8\% & 39.1\% & 65.8\% \\
    G & 0.954 & 18.0\% &  8.9\% & 14.6\% & 35.9\% & 67.4\% \\
    H & 0.987 &  8.7\% &  5.3\% & 12.3\% & 27.8\% & 79.0\% \\
    I & 0.941 & 18.3\% & 11.3\% & 22.8\% & 22.0\% & 58.9\% \\
    J & 0.914 & 24.7\% & 12.0\% & 22.4\% & 27.6\% & 53.0\% \\
    K & 0.954 & 16.2\% &  9.9\% & 21.5\% & 21.3\% & 62.3\% \\
    L & 0.908 & 14.6\% & 21.9\% & 18.5\% & 33.3\% & 66.9\% \\
    M & 0.914 & 19.6\% & 15.1\% & 19.2\% & 26.2\% & 61.2\% \\
    N & 0.928 & 21.9\% & 11.5\% & 20.3\% & 23.6\% & 57.8\% \\
    O & 0.967 & 12.8\% &  8.9\% & 13.0\% & 40.4\% & 74.2\% \\
    P & 0.941 & 12.8\% & 16.1\% & 15.5\% & 29.4\% & 71.7\% \\
    \midrule
    XGB & 0.941 & 28.1\% & 7.3\% & 24.2\% & 24.5\% & 47.7\% \\
    \bottomrule
  \end{tabular}
\end{table}

About a fifth of the nodules receive the single label ``benign'' (more than a quarter with
XGBoost), and roughly 60\% are referred. Two observations follow.

First, the cancers among the auto-benign calls make up 5--22\% of that group: the guarantee limits
how many cancers are missed overall, not how safe a single benign call is, so it does not replace
$R_1$.

Second, the share of cancers covered is below the 0.95 target in 12 of the 16 runs, which looks
like a systematic failure. All runs, however, share one calibration set and one test set, so this
is one draw of the data seen many times. To separate chance from a real shortfall, the model was
held fixed and the 962 calibration and test nodules were randomly re-split 4,000 times, with the
same sizes and class proportions, recalibrating each time (\texttt{experiments/conformal\_checks.py};
Table~\ref{tab:results-resplit}).

\begin{table}[ht]
  \centering
  \small
  \caption{Mondrian split conformal prediction over 4,000 random re-splits of calibration and
           test (model fixed). Mean and standard deviation of the cancer coverage; mean share of
           auto-benign nodules and of cancers among them, with the 5th--95th percentile range.}
  \label{tab:results-resplit}
  \setlength{\tabcolsep}{4pt}
  \begin{tabular}{lrrrr}
    \toprule
    model & \multicolumn{1}{p{1.8cm}}{\raggedleft coverage, our split}
          & \multicolumn{1}{p{2.1cm}}{\raggedleft coverage, re-splits}
          & auto-benign
          & \multicolumn{1}{p{2.5cm}}{\raggedleft cancers in auto-benign} \\
    \midrule
    XGBoost (22 features)        & 0.941 & $0.951 \pm 0.024$ & 25.7\% & 6.4\% (2.3--10.5\%) \\
    LLM, mean of 3 seeds         & 0.934 & $0.950 \pm 0.024$ & 16.7\% & 9.8\% (4.7--14.6\%) \\
    run I                        & 0.941 & $0.951 \pm 0.024$ & 16.9\% & 9.6\% (4.1--14.0\%) \\
    run D                        & 0.908 & $0.951 \pm 0.023$ & 17.6\% & 9.3\% (3.6--14.9\%) \\
    \bottomrule
  \end{tabular}
\end{table}

The average coverage is 0.950--0.951 for every model, exactly as the theory promises, and falls
below 0.95 in about 45\% of the re-splits. Our split was unlucky by about one standard deviation;
the guarantee holds. The same analysis gives the expected share of cancers among the auto-benign
calls: about 10\% for the language model and 6\% for XGBoost, against 35\% in the data and the 5\%
asked of $R_1$.

\chapter{Discussion}
\label{ch:discussion}

\section{Why no decision is certified}
\label{sec:disc-ltt}

The triage layer works as designed, and its statistical code reproduces the worked example of the
design exactly. Two quantities decide whether an auto-benign group can be certified at 5\%:

\begin{itemize}
  \item \textbf{Purity:} how few cancers fall among the lowest-scored nodules. This depends on the
        model.
  \item \textbf{Size:} how many nodules that group contains in the calibration set. With 524
        calibration nodules, a 15\% tail has 78 nodules; LTT needs 59 if none is malignant, 93 if
        one is and 124 if two are.
\end{itemize}

No model reaches both. The best language-model tail (run E) has the purity in its lowest 78
nodules but not the size; the best setting averaged over seeds has 7 cancers there. XGBoost with
the TI-RADS-motivated features, the strongest model in this work (test AUC 0.835), has 2 cancers
among its lowest 78 and 3 among its lowest 104, and is still not certified. The requirement is
not too strict in a statistical sense: it is what 5\% with 95\% confidence means on a few hundred
nodules. A cleaner tail, as a model with AUC of about 0.9 or more would give, or a larger
calibration set are the two ways to reach it.

The experiments show a third lever, in the procedure itself. The planned fixed sequence starts at
the strictest candidate, where only a handful of nodules fall, and stops at the first failure.
For the auto-malignant side of XGBoost this blocks a threshold that would pass on its own. A
sequence that starts where the evidence is, with candidates at quantiles of the validation scores
ordered by their validation $p$-values, keeps the guarantee and turns the auto-malignant side from
almost never certified into certified in 99\% of re-splits for the seed-averaged language model
(Section~\ref{sec:results-ltt-design}). On the auto-benign side no design helps: there the limit
is the model.

Lowering the requirement is a fourth lever. The miss-rate variant and Mondrian split conformal
prediction give a guarantee of a different kind, on the share of cancers missed rather than on
each benign call, and both hold: over 4,000 re-splits of calibration and test, the Mondrian
coverage of cancers averages 0.950--0.951, as the theory says. But the group they call benign
still contains about 10\% cancers for the language model and 6\% for XGBoost, against a 35\%
prevalence. In ACR TI-RADS terms, a 5--20\% risk is that of a TR4 nodule, for which FNA is still
recommended above a size threshold~\citep{tessler2017}. These rules are statistically valid, but
they would not justify skipping a biopsy: \emph{a valid guarantee is not the same as a safe
decision}. Which guarantee is clinically meaningful depends on how a rule will be used, and a
guarantee of the weaker kind should be presented with the share of cancers it leaves among its
benign calls.

The errors that block the guarantee are stable. Most wrong benign calls stay wrong under noise
and under feature removal, so asking the model again does not reveal them. One possible reason is
in the labels: they come from biopsy, and the authors of the dataset report that some biopsy
results disagree with the ultrasound appearance~\citep{gong2022acl}. A malignant nodule that looks
benign on ultrasound cannot be recognised from image features alone, by any model. This is a
hypothesis; it has not been tested here.

\section{What the language model adds}
\label{sec:disc-llm}

The language model learns from the prompts, and it improves when the features improve: by about
0.05 AUC with the revised extraction and by about 0.06 with the TI-RADS-motivated features, the
latter confirmed over three matched seeds. It appears to use the size of the numbers, not only
their order (Section~\ref{sec:results-format}), and plain-language descriptions of the values do
not help it.

Against XGBoost on the same features it stays about 0.05 AUC behind, and the fusion experiment
(Section~\ref{sec:results-fusion}) answers the sharper question of whether it knows something
XGBoost does not: it does not. A stacker given both scores puts all its weight on XGBoost, and the
two models rank the nodules almost as similarly as two training runs of the language model rank
them. In this setting the language model is a slower, noisier and weaker route to the ranking a
tabular model finds directly. Nothing in the experiments indicates that it uses the
\emph{names} of the features, which is the motivation of Ra et al.

This conclusion is limited to what the language model was given: the same numbers a tabular model
sees. Where an LLM could add information is in inputs a tabular model cannot use, such as the
text of radiology reports or the clinical history.

Ra et al.\ report that their language model beats all classical baselines. Several points of their
evaluation make this comparison weaker than it looks:

\begin{itemize}
  \item \textbf{Feature names.} Replacing the names with generic labels lowers their internal test
        AUC from 0.700 to 0.695 and their external test AUC from 0.862 to 0.849. These differences
        are called significant, but no test or interval is given.
  \item \textbf{AUC of the baselines.} For all seven scikit-learn baselines, the reported AUC is
        within 0.004 of the accuracy (for example 0.631 and 0.631 for the random forest), while
        for their own model and for TransTab the AUC is clearly higher than the accuracy (0.700
        against 0.671). This pattern is expected if the baselines' AUC was computed from 0/1
        predictions, which turns AUC into balanced accuracy. The comparison would then favour the
        language model.
  \item \textbf{Fine-tuning.} Training only the classification head of the frozen model gives an
        AUC of 0.699, against 0.700 with LoRA; without mRMR, head-only training gives 0.710, the
        highest AUC in their ablation table.
  \item \textbf{Labels.} Their labels are BI-RADS categories (3 benign, 4--6 malignant), that is a
        radiologist's assessment of the image, not a biopsy result.
  \item \textbf{Size and uncertainty.} The internal test set has 235 images and the external one
        about 55 after half of INbreast was used for retraining. No intervals are reported, a
        single run is shown, and the model was trained for over 500 epochs and selected on
        validation. In this work the seed alone moved the test AUC by up to 0.036, more than most
        of the differences in their tables.
\end{itemize}

With their backbone and training settings on TN3K (run G: Llama-2-7B, learning rate $10^{-4}$,
LoRA on all linear layers, 16 features instead of 8), the test AUC is 0.731, below Qwen2.5-1.5B
(0.754) and XGBoost with the same features (0.805). The results of this thesis support a cautious
reading of theirs: a language model can classify radiomics prompts, but below a properly evaluated
tabular model and without adding information to it.

\section{Limitations}
\label{sec:disc-limitations}

\begin{itemize}
  \item \textbf{One 2D image per nodule.} A nodule is a 3D structure seen in several planes and in
        real time; one still image gives a partial view.
  \item \textbf{Ultrasound images.} Brightness depends on the device, its settings and the
        operator. Per-image rescaling removes part of this variation, and with it part of the
        information on absolute echogenicity.
  \item \textbf{Clinical signs approximated.} The TI-RADS-motivated features are heuristics, not
        expert ratings. Echogenicity relative to thyroid tissue and echogenic foci are not
        captured well, and whether a view is transverse is unknown.
  \item \textbf{Patient overlap.} The pooled split is drawn at image level, so the absolute test
        metrics are likely optimistic for new patients (Section~\ref{sec:data-patients}). The
        conformal guarantees remain valid for new images from the same pool, not for new
        patients.
  \item \textbf{Single runs.} Apart from the seed repeats of the main comparison, each setting was
        trained once. The seed alone moves the test AUC by up to about 0.04, so single-run
        differences of that size are not interpreted.
  \item \textbf{Intervals.} The bootstrap intervals hold the trained model fixed and resample
        images, not patients; they omit training noise and patient clustering and are somewhat
        too narrow. Many comparisons were made without a multiplicity correction; the main
        findings hold across seeds, but borderline single intervals should not be over-read.
  \item \textbf{Choices informed by results.} Dropping the seventh TI-RADS feature, keeping the
        decimal format and the validation-ordered LTT were decided after seeing calibration or
        test results. The main conclusions do not depend on these choices.
  \item \textbf{Hyperparameters.} The learning rate was chosen for the 1.5B model. The 7B Qwen
        model was tried at two learning rates, and Llama-2-7B used the settings of Ra et al.\ but
        5 epochs instead of 500. The large-model runs used one seed each, and the validation AUC of
        the low-learning-rate 7B run was still rising at the last epoch, so a tuned 7B model was
        not tested. XGBoost was not tuned.
  \item \textbf{Prevalence.} With 35\% cancers, TN3K is a referral (FNA) population. Risks
        defined among the automatic decisions, such as $R_1$, depend on prevalence and would be
        easier to meet, and different, in a screening population.
  \item \textbf{One dataset.} All results come from TN3K.
  \item \textbf{Exploratory analyses.} The miss-rate variant, the validation-ordered LTT and the
        re-split analysis were added after the main results.
\end{itemize}

\chapter{Ongoing and future work}
\label{ch:future}

The questions about the classifier are answered: the best setting, its training noise, the effect
of the TI-RADS-motivated features, model size, prompt format and the comparison with XGBoost. The
open questions concern the triage layer and the generality of the results.

\section{A confirmatory Learn-then-Test design}
\label{sec:future-running}

The planned fixed sequence stopped at candidates too small to pass
(Section~\ref{sec:results-ltt}). A design fixed in advance can avoid this without giving up the
guarantee: candidate thresholds at quantiles of the validation scores, so that every candidate
holds enough calibration nodules, tested in the order of their validation $p$-values. Run on new
random splits, it would show whether the exploratory auto-malignant certification of
Section~\ref{sec:results-ltt} is reproducible, and on which side a guarantee is within reach of
the current models.

\section{Conformal results over several splits}

All conformal numbers come from one calibration set of 524 nodules. The re-split analysis
(Table~\ref{tab:results-resplit}) varies the calibration and test assignment for a fixed model;
retraining on several random splits would add the variation of the model itself and would give
the expected share of cancers among the automatic benign calls with an honest range.

\section{Prompt variants not yet run}
\label{sec:future-prompts}

\begin{itemize}
  \item \textbf{Meaning without numbers.} The \emph{semantic only} format, which gives the band and
        the phrase without any number, is implemented. Since the model uses the magnitude of the
        values, it is expected to lose; it would show how much can be learnt from meaning alone.
  \item \textbf{TI-RADS features alone.} Phrases such as ``taller than wide'' or ``a more lobulated
        outline'' are the language of radiology reports. A prompt with the six TI-RADS-motivated
        features only would test whether meaning helps where the feature names mean something
        to a language model.
\end{itemize}

\section{A second dataset}

ThyroidXL will be used as a separate dataset, not merged with TN3K. It can show whether the gain
from the TI-RADS-motivated features and the gap between the language model and XGBoost hold on
other data and, with more calibration nodules, whether the triage guarantee becomes reachable.

\section{Inputs a tabular model cannot use}

The fusion experiment shows that the language model adds nothing when it sees the same numbers as
XGBoost. The case for a language model rests on inputs that a tabular model cannot read, such as
free-text radiology reports or clinical history. Combining them with the radiomics features is the
setting in which a language model could contribute information of its own.


\backmatter
\cleardoublepage
\phantomsection
\addcontentsline{toc}{chapter}{\bibname}
\bibliographystyle{plainnat}
\bibliography{references}

\end{document}
